\documentclass[10pt,a4paper]{article}

\usepackage[T1]{fontenc}
\usepackage[utf8]{inputenc}
\usepackage[ngerman]{babel}
\usepackage[a4paper,margin=12mm]{geometry}
\usepackage{xcolor}
\usepackage{lmodern}
\usepackage{microtype}
\usepackage{listings}
\usepackage{tikz}
\usetikzlibrary{shapes.multipart,positioning}

\definecolor{Navy}{HTML}{17324D}
\definecolor{Blue}{HTML}{3B6EA5}
\definecolor{LightBlue}{HTML}{EAF3FA}
\definecolor{Mint}{HTML}{E8F5EF}
\definecolor{Line}{HTML}{D8E1E8}
\definecolor{Muted}{HTML}{617080}

\setlength{\parindent}{0pt}
\setlength{\parskip}{2pt}
\pagestyle{empty}

\lstset{
  language=Java,
  basicstyle=\ttfamily\footnotesize,
  keywordstyle=\color{Blue}\bfseries,
  commentstyle=\color{Muted},
  backgroundcolor=\color{LightBlue},
  frame=single,
  rulecolor=\color{Line},
  columns=fullflexible,
  keepspaces=true,
  showstringspaces=false,
  aboveskip=3pt,
  belowskip=3pt
}

\newcommand{\tasktitle}[1]{%
  \vspace{2pt}{\large\bfseries\color{Blue}#1}\par
  \color{Line}\hrule\color{black}\vspace{2pt}
}

\begin{document}

\colorbox{Navy}{%
  \parbox{\dimexpr\textwidth-2\fboxsep\relax}{%
    \color{white}
    {\small\bfseries INFORMATIK · KLASSE 10\hfill LÖSUNGEN}\par
    \vspace{4pt}
    {\LARGE\bfseries Sicherungsblatt: Aufgabe 2}\par
    {\large Wiederholung Roboter}
  }%
}

\vspace{4pt}
\colorbox{LightBlue}{%
  \parbox{\dimexpr\textwidth-2\fboxsep\relax}{%
    \textbf{\color{Navy}Grundidee:}
    Der Roboter wiederholt Anweisungen, solange vor ihm keine Wand ist.
    An einer Ecke dreht er sich. Für eine vier Ziegel hohe Mauer darf er
    nicht auf die Mauer steigen, sondern muss auf einer Spur daneben laufen.
  }%
}

\tasktitle{a) Ziegelreihe bis zur Wand}

Eine mögliche Schleife lautet:

\begin{lstlisting}
Robot roboter = new Robot(1, 1, 15, 15);

while (!roboter.istWand()) {
   roboter.hinlegen();
   roboter.schritt();
}
\end{lstlisting}

\textbf{Merksatz:} Eine \texttt{while}-Schleife eignet sich, wenn vor dem
Start nicht feststeht, wie oft wiederholt werden muss.

\tasktitle{b) Ziegel am gesamten Rand entlang}

Eine Seite und das anschließende Drehen werden viermal wiederholt:

\begin{lstlisting}
Robot roboter = new Robot(1, 1, 15, 15);

for (int seite = 0; seite < 4; seite++) {
   while (!roboter.istWand()) {
      roboter.hinlegen();
      roboter.schritt();
   }
   roboter.linksDrehen();
}
\end{lstlisting}

Je nach Startpunkt und gewünschter Lage der Ziegel kann eine zusätzliche
Drehung oder ein versetzter Startpunkt nötig sein. Entscheidend ist die
Schachtelung: Die innere Schleife bearbeitet eine Seite, die äußere alle vier.

\tasktitle{c) Mauer mit Höhe 4}

\texttt{hinlegen(4)} legt vier Ziegel auf das Feld vor dem Roboter. Weil der
Roboter höchstens einen Ziegel hochspringen kann, läuft er auf einer Spur
\emph{neben} der Mauer. Das Grundmuster für eine gerade Seite ist:

\begin{lstlisting}
while (!roboter.istWand()) {
   roboter.hinlegen(4);

   // Zur Laufrichtung drehen, ein Feld weitergehen
   // und wieder zur Mauer drehen.
   roboter.linksDrehen();
   roboter.schritt();
   roboter.rechtsDrehen();
}
\end{lstlisting}

Für jede Ecke müssen Laufrichtung und Abstand zum Rand angepasst werden.
Eine korrekte Lösung darf Hilfsmethoden verwenden, zum Beispiel
\texttt{baueSeite()} und \texttt{geheZurNaechstenSeite()}.

\tasktitle{d) Struktogramm zu c)}

\begin{center}
\begin{tikzpicture}[
  node distance=0pt,
  every node/.style={
    draw=Navy,
    minimum width=11.8cm,
    text width=11.2cm,
    align=left,
    font=\small,
    inner sep=5pt
  }
]
  \node (start) {\textbf{Roboter erzeugen und an der inneren Randspur ausrichten}};
  \node[below=of start] (outer) {\textbf{Wiederhole für jede Seite:}};
  \node[below=of outer, fill=LightBlue] (inner) {
    \textbf{Solange das nächste Feld an der Mauer bearbeitet werden kann:}\\
    vier Ziegel ablegen; zur Laufrichtung drehen; einen Schritt gehen;
    wieder zur Mauer drehen
  };
  \node[below=of inner] (corner) {\textbf{Ecke:} zur nächsten Randseite wechseln und neu ausrichten};
\end{tikzpicture}
\end{center}

\vspace{3pt}
\colorbox{Mint}{%
  \parbox{\dimexpr\textwidth-2\fboxsep\relax}{%
    \textbf{\color{Navy}Prüfkriterien:}
    Das Programm endet ohne Wandkollision, bearbeitet alle gewünschten
    Randfelder, erzeugt überall genau vier Ziegel und das Struktogramm zeigt
    dieselben Schleifen, Bedingungen und Anweisungen wie der Programmcode.
  }%
}

\vspace{3pt}
{\small\color{Muted}
Informatik 10 · Modellierung und Programmierung · Aufgabe 2}

\end{document}
